\documentclass[a4paper,10pt]{article}
\usepackage{geometry}
\usepackage{enumitem}
\usepackage[hidelinks]{hyperref}
\usepackage{titlesec}
\usepackage{xcolor}
\geometry{margin=0.8in}
\pagenumbering{gobble}
\titleformat{\section}{\large\bfseries\color{blue!70!black}}{}{0em}{}[\titlerule]
\titlespacing*{\section}{0pt}{2pt}{2pt}
\setlist[itemize]{noitemsep, topsep=2pt, leftmargin=1.5em}
\begin{document}

\begin{center}
\Large \textbf{Zhalok Rahman}\\[3pt]
\normalsize Software Engineer 2 at Datatree Technology\\[2pt]
\small rahmanzhalok@gmail.com \\phone +8801716922067\\[2pt]
\small LinkedIn: https://www.linkedin.com/in/zhalok-rahman/\\[2pt]
\small Github: https://github.com/zhalok\\[2pt]
\small Video Introduction: https://www.loom.com/share/c5ab46a4737244edafd9f052555767a1\\[2pt]
\end{center}
\section*{Skills}
Golang, Python (FastAPI, Asyncio, Pytorch, Pandas), Node JS (Nest JS, Express JS, Koa), JavaScript/TypeScript (React JS, Next JS, PostgreSQL, MongoDB, Redis (Cache, Pub/Sub, Streams), RabbitMQ, WebSocket, SSE, Docker, AWS (ECS, CloudWatch, SQS, Kinesis, Lambda, OpenSearch, DynamoDB), Grafana, Prometheus, CI/CD (Github Actions, CodePipeline
\vspace{2pt}
\section*{Experience}
\textbf{Kinetik Healthcare} \hfill \textit{Software Engineer 2}\\
\textit{2025-04-01 -- Present}
\begin{itemize}
\item Implemented push based realtime system using Redis PubSub
\item Generated notifications and pushed in Redis Streams via MongoDB CDC realm trigger
\item Implemented RBAC validation and realtime access revoke on role change
\item Built endpoints for AI agent tool calls
\item Built AI call center agent with ElevenLabs
\item Integrated human call center as fallback for AI callcenter service
\item Managed ElevenLabs AI agent as code with \textbf{Python} SDK, building a CI/CD deploy pipeline that fetches the agent config from the development agent, replicates it into staging/production with injected environment config, and avoids shipping UI-edited agents directly for better auditability
\item Built a \textbf{Python} voice service (ai-service) with HTTP and WebSocket endpoints: Twilio hits the HTTP endpoint to start a session, receives TwiML pointing to the WebSocket endpoint with the session id, then streams voice over that WebSocket connection
\item Streamed incoming voice audio from the WebSocket connection to the LLM orchestration service over gRPC streams
\item Integrated Langfuse for tracing and monitoring the voice AI agent pipeline
\item Managed infrastructure in AWS ECS
\item Debugged build and deployment issues
\end{itemize}
\textbf{PropertyFinder} \hfill \textit{Software Engineer}\\
\textit{2024-10-01 -- 2026-03-31}
\begin{itemize}
\item Worked on B2B modernization with Golang and AWS based services like DynamoDB, Lambda, SQS, ECS
\item Performed load testing experiment with go script and created a dashboard with necessary metrics to observe the changes in high load
\item Implemented asynchronous business logics with SQS, Lambda (go) and Kinesis streams
\item Fixed and investigated production incidents
\item Collaborated with data engineering team for fixing data discrepancy issues during migration
\item Written go scripts for resolving incident and migration of data
\item Contributed to data migration using \textbf{Python} and Apache Spark scripts
\item Fixed github actions deployment failures
\end{itemize}
\textbf{Learnyx AI} \hfill \textit{Fullstack Engineer}\\
\textit{2023-10-01 -- 2025-02-01}
\begin{itemize}
\item Built an ai-service entirely in \textbf{Python}, consuming RabbitMQ as the entrypoint invocation for heavy AI workloads like vector search and LLM orchestration
\item Built AI generation service with Langchain, FastAPI and OpenAI models
\item Built websocket and server sent events endpoints for realtime AI generation delivery
\item Created RabbitMQ consumers for handling heavy workloads asyncly with PikaClient
\item Offloaded event loop blocking sync tasks into background tasks
\item Scheduled gathered coroutines for concurrent resolutions
\item Traced event loop blocking issues with sync operations and optimized the response time
\end{itemize}
\textbf{Innovext IT} \hfill \textit{Fullstack Developer}\\
\textit{2022-12-01 -- 2023-10-01}
\begin{itemize}
\item Built REST apis with Nest JS, Express JS and FastAPI (for deploying  fine-tuned model)
\item Implemented payment and monetary transaction features by integrating third party payment gateways (stripe and paypal)
\item Built client side application with Next JS, used Redux for efficient data fetching and dataflow
\item Built an internal monetization system with credits, ensured smooth credit transfer with mongodb sessions
\item Built server with FastAPI for inferencing fine-tuned object detection model in onnx format
\item Created scripts for data preprocessing and training models
\end{itemize}
\vspace{2pt}
\section*{Projects}
\textbf{Remote Browser Control}\\
\textit{Remote-controlled Playwright browser session streamed to a React client over WebSocket.}
\begin{itemize}
\item Built FastAPI app for receiving commands for playwright to apply on browser session via websocket and emit current snapshot of the browser
\item Created a front end app with react to render the remote browser session and also send commands via websocket
\item Demo: https://bit.ly/4wDD7om
\item Github: https://github.com/zhalok/remote-browser-control
\end{itemize}
\small Tech: \small{Playwright}, \small{Python}, \small{FastAPI}, \small{WebSocket}\\
\textbf{Invoice Extractor}\\
\textit{Automated reading of supplier invoices and verifying the extracted data before registering into a (mock) accounting system}
\textit{https://github.com/zhalok/invoice-extractor.git}
\begin{itemize}
\item Built a FastAPI upload/job-polling API and a \textbf{Python} RabbitMQ worker pool running an extract, hydrate, verify, and register pipeline
\item Designed dual auto-submit and manual-review flows: verified invoices auto-register immediately, while unresolved supplier, duplicate, or math-mismatch cases fall back to a human review UI before submission
\item Github: https://github.com/zhalok/invoice-extractor.git
\item Demo: https://youtu.be/BsikeXZiBuI
\end{itemize}
\small Tech: \small{**Python**}, \small{FastAPI}, \small{RabbitMQ}\\
\vspace{2pt}
\section*{Publications}
\textbf{Optical Mark Recognition with Object Detection and Clustering}\\
\small 2024 6th International Conference on Electrical Engineering and Information and Communication Technology (ICEEICT)
\hfill \textit{May 2024}\\
\small \url{https://doi.org/10.1109/ICEEICT62016.2024.10534561}\\
\begin{itemize}
\item Written \textbf{Python} scripts for generating synthetic training data
\item Built training scripts in \textbf{Python} with ultralytics for training lightweight YOLO models
\item Implemented clustering based unsupervised algorithm for column detection and serialization of answers
\item Exported finetuned model in Onnx format for inferencing via API
\item demo: https://www.youtube.com/watch?v=BiEybE5qQqw
\end{itemize}
\vspace{2pt}
\section*{Certifications}
\textbf{AWS Certified Cloud Practitioner}\\
\small AWS\\
\small https://www.credly.com/badges/201260d3-12e0-4d22-a49f-4e64d68b348e\\
\textit{Deep dive on AWS services}\\
\textbf{Optimizing MongoDB Performance with Tuning Tools}\\
\small MongoDB\\
\small https://www.credly.com/badges/e861b7b4-028a-47d4-aed6-7e586d4d95da\\
\textit{Learned about diagnosing and tuning MongoDB performance using database tooling}\\
\textbf{Supervised Machine Learning: Regression and Classification}\\
\small Coursera\\
\small https://coursera.org/share/0ef4defbf44ee7d868c52c7e1a3bb94f\\
\textit{Supervised learning from scratch, learning about mathematical algorithms working as backbone of machine learning models}\\
\vspace{2pt}
\section*{Education}
\textbf{Shahjalal University of Science and Technology}\\
Bachelors of Science in Computer Science\\
\textit{2019-01-01 -- 2024-04-01}\\
\vspace{2pt}
\end{document}